\section{Conjuntos de datos de post-entrenamiento: tres dimensiones de la calidad}

\subsection{Definición de buenos datos}

Quien imparte la clase considera que la calidad de los datos es el problema más central de todo el proceso de post-entrenamiento, y puede medirse a partir de tres dimensiones:

\begin{importantbox}{Marco tridimensional de la calidad de los datos}
\begin{enumerate}
    \item \textbf{Accuracy (exactitud)}: si la respuesta es correcta. Es relativamente fácil de verificar para preguntas simples; para código puede verificarse mediante unit tests, y para matemáticas mediante un solver.
    \item \textbf{Diversity (diversidad)}: si el conjunto de ejemplos cubre los distintos escenarios de interacción que puede tener el usuario. Una diversidad insuficiente provoca que al modelo le falte gravemente capacidad en ciertos dominios.
    \item \textbf{Complexity (complejidad)}: si los ejemplos son lo bastante desafiantes, con razonamiento de varios pasos, chain-of-thought, etc.
\end{enumerate}
\end{importantbox}

\subsection{La importancia de la diversidad}

Al visualizar en 2D el embedding de distintos conjuntos de datos, la diferencia se aprecia de manera directa:

\begin{itemize}
    \item \textbf{Conversaciones reales de usuarios} (como ShareGPT, WildChat): la distribución del embedding cubre un rango sumamente amplio, porque las preguntas de los usuarios abarcan todo tipo de temas
    \item \textbf{Conjuntos de datos matemáticos}: el embedding se concentra en gran medida en una sola región
    \item \textbf{Conjuntos de datos de código}: también se concentran en una región limitada
\end{itemize}

Por lo tanto, si el objetivo es construir un LLM de propósito general, los datos de matemáticas y de código por sí solos son claramente insuficientes; es necesario incorporar datos de conversación diversos para ampliar la cobertura.

\begin{warningbox}{La trampa de la diversidad en los datos sintéticos}
Uno de los grandes riesgos de los datos sintéticos (synthetic data) es el colapso de la diversidad. Una pequeña cantidad de datos sintéticos puede ampliar la diversidad de manera efectiva, pero un \textbf{exceso} de datos sintéticos, por el contrario, provoca que la distribución de los datos se estreche y que la diversidad caiga de forma abrupta, lo que termina por reducir la calidad general de los datos.
\end{warningbox}

\subsection{Aumento de la complejidad: Auto-Evol}

Un ejemplo de baja complejidad como «¿Qué tan alta es la Torre Eiffel?» puede responderse con solo unas cuantas palabras. Es posible usar marcos como \textbf{Auto-Evol} para que el LLM reescriba el prompt y aumente su complejidad; por ejemplo, exigiendo un análisis desde varios ángulos, comparando datos de distintas épocas, explicando principios de ingeniería, etc. Las preguntas reescritas requieren respuestas más largas y detalladas, lo que aumenta de manera natural la complejidad de los datos.

\subsection{Formato de los datos}

El post-entrenamiento utiliza dos formatos de datos principales:

\begin{table}[H]
\centering
\caption{Formatos de datos de post-entrenamiento}
\begin{tabular}{lll}
\toprule
\textbf{Formato} & \textbf{Composición} & \textbf{Uso} \\
\midrule
Instruction Data & system prompt + instruction + output & SFT \\
Preference Data & instruction + chosen answer + rejected answer & DPO/PPO \\
\bottomrule
\end{tabular}
\end{table}

\subsection{Resumen de la sección}

La calidad de los datos está determinada en conjunto por tres dimensiones: exactitud, diversidad y complejidad. Al construir un conjunto de datos, hay que prestar especial atención al riesgo de diversidad de los datos sintéticos, y pueden usarse herramientas como Auto-Evol para aumentar la complejidad de los ejemplos.

\newpage
